\begin{question}

  \begin{questionpart}
  Show that the determinant of a $2 \times 2$ matrix
    $\vec{\sigma} \cdot \vec{a}$ is invariant under

    \begin{equation*}
      \vec{\sigma} \cdot \vec{a} \rightarrow \vec{\sigma} \cdot \vec{a}' \equiv
      exp \left( \frac{i\vec{\sigma} \cdot \hat{n}\phi}{2}\right)
      \vec{\sigma} \cdot \vec{a} \ 
      exp \left( \frac{-i\vec{\sigma} \cdot \hat{n}\phi}{2}\right)
    \end{equation*}

    where the matrices $\sigma$ are given in (3.50).\\
  \end{questionpart}

  \begin{questionpart}
  Find $a_k'$ in terms of $a_k$ when $\hat{n}$ is in the
    positive z-direction and interpret your result.\\
  \end{questionpart}

\end{question}

\begin{answer}

%%%%%%%%%%%%%%%%%%%%%%%%%%%%%%%%%%%%%%%%%%%%%%%%%%%%%%%%%%%%%%%%%%%%%%%%%%%%%%%
%%                              Part a                                       %%
%%%%%%%%%%%%%%%%%%%%%%%%%%%%%%%%%%%%%%%%%%%%%%%%%%%%%%%%%%%%%%%%%%%%%%%%%%%%%%%

\begin{answerpart}{Equivalence of determinants}

  We start by separating out the three terms individually
  
  \begin{equation}
    \begin{split}
      X
      &= \vec{\sigma} \cdot \vec{a} \rightarrow \vec{\sigma} \cdot \vec{a}'\\
      &\equiv
      exp \left( \frac{i\vec{\sigma} \cdot \hat{n}\phi}{2}\right)
      \vec{\sigma} \cdot \vec{a} \ 
      exp \left( \frac{-i\vec{\sigma} \cdot \hat{n}\phi}{2}\right)\\
      &= A \cdot B \cdot C\\
    \end{split}
  \end{equation}

  Using the separability of the determinant of the factors:

  \begin{equation}
    det(X) = det(ABC) = det(A) det(B) det(C)
  \end{equation}

  Ultimately, we are trying to show that:

  \begin{equation}
    \begin{split}
      det(X)
      &= det(B) \\
      &= det(ABC) \\
      &= det(A) det(B) det(C) \\
    \end{split}
  \end{equation}

  This implies that
  
  \begin{equation}
    \begin{split}
      \frac{det(X)}{det(B)}
      &= 1 \\
      &= det(A)det(C) \\
      &= det(AC) \\
    \end{split}
  \end{equation}

  Putting it all together, we see that:
  
  \begin{equation}
    \begin{split}
      det(X)
      &= det \left\{ exp \left( \frac{i\vec{\sigma} \cdot \hat{n}\phi}{2}\right)
      \vec{\sigma} \cdot \vec{a} \ 
      exp \left( \frac{-i\vec{\sigma} \cdot \hat{n}\phi}{2}\right) \right\}\\
      &= det(ABC) \\
      &= det(A)det(B)det(C) \\
      &= det(A)det(C)det(B) \\
      &= det(AC)det(B) \\
      &= det(I)det(B) \\
      &= det(B) \\
      &= det(\vec{\sigma} \cdot \vec{a})
      \qed\\
    \end{split}
  \end{equation}

\end{answerpart}


%%%%%%%%%%%%%%%%%%%%%%%%%%%%%%%%%%%%%%%%%%%%%%%%%%%%%%%%%%%%%%%%%%%%%%%%%%%%%%%
%%                              Part b                                       %%
%%%%%%%%%%%%%%%%%%%%%%%%%%%%%%%%%%%%%%%%%%%%%%%%%%%%%%%%%%%%%%%%%%%%%%%%%%%%%%%

\begin{answerpart}{\texorpdfstring{$\hat{n} = \hat{z}$}{n = z}}

To address this part, we must first examine the meaning of
$\vec{\sigma} \cdot \hat{n}$.  The $n$ vector is a vector of real
numbers between -1 and 1 indicating the fraction within $\sigma_i$.
The problem specifies $\hat{n}$ = $\hat{z}$ meaning that

\begin{equation}
  \sigma \cdot \hat{n} = \sigma_z
\end{equation}

Next, we'll substitute that back into our definition of $a'$

\begin{equation}
  \inlabel{b-characteristic}
  \begin{split}
    A &=
    \vec{\sigma} \cdot \vec{a}' \\
    &= exp \left( \frac{i\sigma_z\phi}{2}\right)
       \vec{\sigma} \cdot \vec{a} \ 
       exp \left( \frac{-i\sigma_z\phi}{2}\right)\\
  \end{split}
\end{equation}

To solve for $A$, we'll need to use the definition of matrix
exponentiation:

\begin{equation}
  exp(X) \equiv \sum_{0}^{\infty} \frac{1}{k!} X^k
\end{equation}

Applying this definition to $A$, we find that:

\begin{equation}
  \begin{split}
    A &=
    \vec{\sigma} \cdot \vec{a}' \\
    &= exp \left( \frac{i\sigma_z\phi}{2}\right)
       \vec{\sigma} \cdot \vec{a} \ 
       exp \left( \frac{-i\sigma_z\phi}{2}\right)\\
    &= \sum_{0}^{\infty}{\frac{1}{k!}\left(\frac{i\sigma_z\phi}{2}\right)^k}
       \vec{\sigma} \cdot \vec{a} \
       \sum_{0}^{\infty}{\frac{1}{k!}\left(\frac{-i\sigma_z\phi}{2}\right)^k}
  \end{split}
\end{equation}

Because $\sigma_i^2 = I$ and owing to the nature of the i in the
exponent, we can see that the sum is going to look interesting, so
let's explore that.

\begin{equation}
  \begin{array}{ccccc}
    Term & Imag & Sign & \sigma \\
    \hline
    0 & 1 & + & I \\
    1 & i & + & \sigma_z \\
    2 & 1 & - & I \\
    3 & i & - & \sigma_z \\
    \hline
    4 & 1 & + & I \\
    5 & i & + & \sigma_z \\
    6 & 1 & - & I \\
    7 & i & - & \sigma_z \\
    \hline
    ...
  \end{array}
\end{equation}

The periodicity of 4 and odd/even behavior allows us to rewrite the
original sum as follows:

\begin{equation}
  \sum_{0}^{\infty}{\frac{1}{k!}\left(\frac{i\sigma_z\phi}{2}\right)^k}
  = \sum_{0}^{\infty}{\frac{(-1)^k}{2k!}\left(\frac{\phi}{2}\right)^{2k}}
  + i\sigma_z\sum_{0}^{\infty}{\frac{(-1)^k}{(2k+1)!}\left(\frac{\phi}{2}\right)^{2k+1}}
\end{equation}

As it happens, this pattern matches Euler's Taylor series formulation
of $cos + i\ sin$ meaning that:

\begin{equation}
  \sum_{0}^{\infty}{\frac{1}{k!}\left(\frac{i\sigma_z\phi}{2}\right)^k}
  = \cos(\frac \phi 2) + i \sigma_z \sin(\frac \phi 2)
\end{equation}

Substituting that back into our initial equation (\inref{b-characteristic})

\begin{equation}
  \begin{split}
    A &=
    \vec{\sigma} \cdot \vec{a}' \\
    &= exp \left( \frac{i\sigma_z\phi}{2}\right)
       \vec{\sigma} \cdot \vec{a} \ 
       exp \left( \frac{-i\sigma_z\phi}{2}\right)\\
    &= \left\{ \cos(\frac \phi 2) + i \sigma_z \sin(\frac \phi 2) \right\}
       \vec{\sigma} \cdot \vec{a} \ 
       \left\{ \cos(-\frac \phi 2) + i \sigma_z \sin(-\frac \phi 2) \right\}\\
    &= \left\{ \cos(\frac \phi 2) + i \sigma_z \sin(\frac \phi 2) \right\}
       \vec{\sigma} \cdot \vec{a} \ 
       \left\{ \cos(\frac \phi 2) - i \sigma_z \sin(\frac \phi 2) \right\}\\
  \end{split}
\end{equation}

With the help of a couple of identities:

\begin{equation}
  \cos(\frac \phi 2) + i\sin(\frac \phi 2) = e^{i \tfrac \phi 2}
\end{equation}

\begin{equation}
  \begin{split}
    e^{-i \tfrac \phi 2}
    &= (\cos(-\frac \phi 2) + i\sin(-\frac \phi 2) \\
    &= \cos(\frac \phi 2) - i\sin(\frac \phi 2) \\
  \end{split}
\end{equation}

  And recalling equation 3.50 from the book:  
  
  \begin{equation*}
    \tag{3.50}
    \sigma_1 =
    \begin{pmatrix}
      0 & 1 \\
      1 & 0 \\
    \end{pmatrix},
    \sigma_2 =
    \begin{pmatrix}
      0 & -i \\
      i &  0 \\
    \end{pmatrix},
    \sigma_3 =
    \begin{pmatrix}
      1 &  0 \\
      0 & -1 \\
    \end{pmatrix}
  \end{equation*}

  We can rewrite our matrix as

  \begin{equation}
    \begin{split}
      A &=
      \begin{pmatrix}
        e^{i \tfrac \phi 2} & 0\\
        0 & e^{-i \tfrac \phi 2}\\
      \end{pmatrix}
      \vec{\sigma} \cdot \vec{a}
      \begin{pmatrix}
        e^{-i \tfrac \phi 2} & 0\\
        0 & e^{i \tfrac \phi 2}\\
      \end{pmatrix}\\
    \end{split}
  \end{equation}

Recalling equation 3.56:

\begin{equation*}
  \tag{3.56}
  \begin{split}
    \vec{\sigma} \cdot \vec{a} &\equiv \sum_k a_k \sigma_k \\
    &= \begin{pmatrix}
      a_3        & a_1 - ia_2 \\
      a_1 + ia_2 & -a_3 \\
    \end{pmatrix}
  \end{split}
\end{equation*}

  We can now construct the final values for $\vec{\sigma} \cdot
  \vec{a'}$:

  \begin{equation}
    \inlabel{final}
    \begin{split}
      \vec{\sigma} \cdot \vec{a'}
      &=
      \begin{pmatrix}
        e^{i \tfrac \phi 2} & 0\\
        0 & e^{-i \tfrac \phi 2}\\
      \end{pmatrix}
      \begin{pmatrix}
        a_3  & a_1 - ia_2 \\
        a_1 + ia_2 & -a_3 \\
      \end{pmatrix}
      \begin{pmatrix}
        e^{-i \tfrac \phi 2} & 0\\
        0 & e^{i \tfrac \phi 2}\\
      \end{pmatrix}\\
      &=
      \begin{pmatrix}
        e^{i \tfrac \phi 2} & 0\\
        0 & e^{-i \tfrac \phi 2}\\
      \end{pmatrix}
      \begin{pmatrix}
        e^{-i \tfrac \phi 2}(a_3)  & e^{i \tfrac \phi 2}(a_1 - ia_2) \\
        e^{-i \tfrac \phi 2}(a_1 + ia_2) & e^{i \tfrac \phi 2}(-a_3) \\
      \end{pmatrix}\\
      &=
      \begin{pmatrix}
        a_3  & e^{i\phi}(a_1 - ia_2) \\
        e^{-i\phi}(a_1 + ia_2) & -a_3 \\
      \end{pmatrix}\\
    \end{split}
  \end{equation}

What is truly fascinating here is that the rotation has been applied
to the anti-diagonal components of the final matrix.  If a rotation
about the z-axis were to occur, then the z-axis components would not
change.  Sakurai's selection of the z-axis for this rotation is
presumably intended to help illustrate this point.  Instead of
modifying the z-axis components, we instead see that the x and y
components are rotated.

Comparing 3.56 to \inref{final}, we see that while $a_3' = a_3$, terms
containing $a_1$ and $a_2$ each have a factor of $e^i\tfrac \phi 2$
applied.
\end{answerpart}
\end{answer}
